\section{Preliminaries}


\subsection{Varieties and rational maps}

  Throughout this note, all schemes are of finite type and separated over a field or a noetherian integral scheme.
  A variety is a geometrically integral scheme of finite type and separated over a field.
  In other words, when base field is algebraically closed, a variety is an integral scheme of finite type and separated over the base field.
  And in general case, $X$ is a variety if and only if its base change to the algebraic closure of the base field is a variety.
  We usually use $\kk$ to denote an arbitrary field and $\kkk$ to denote an algebraically closed field.

  When work over a non-algebraically closed field $\kk$, we often consider the set $X(\kkk)$ of $\kkk$-points of a variety $X$ over $\kk$ rather than the scheme-theoretic points of $X$.
  Recall that for any field extension $\kk'/\kk$, the set of $\kk'$-points of $X$ is given by
  \[ X(\kk') = \Mor_{\kk}(\Spec \kk', X) = \{ x: \Spec \kk' \to X \text{ over } \kk \}. \]
  A point $x \in X(\kkk)$ is determined by the image of $x$ in $X$ and the induced embedding $\kappa(x) \hookrightarrow \kk'$ of the residue field $\kappa(x)$ of $x$ into $\kk'$.
  For a scheme-theoretic closed point $\xi \in X$, we will say that $\xi \in X(\kk')$ if the residue field $\kappa(\xi)$ of $\xi$ is a subfield of $\kk'$, or equivalently, there exists $x \in X(\kk')$ such that the image of $x$ in $X$ is $\xi$.
  Note that there may be at most $[\kappa(\xi):\kk]$ different points in $X(\kk')$ with the same image $\xi$ in $X$.
  \begin{proposition}\label{prop:closed_points_on_varieties_over_non-algebraically_closed_fields}
    Let $X,Y$ be varieties over $\kk$ and $f: X \to Y$ be a morphism.
    Then 
    \begin{enumerate}
      \item $X(\kkk)$ is canonically identified with the set of closed points of $X_\kkk$;
      \item $f: X(\kkk) \to Y(\kkk), x \mapsto f(x) = x \circ f$ coincides with the morphism $f_\kkk: X_\kkk \to Y_\kkk$ on the closed points under the above identification.
    \end{enumerate}
  \end{proposition}
  \begin{remark}
    If we consider the set of scheme-theoretic points of $X$ rather than $\kkk$-points, then the morphism $f: X \to Y$ may behave unlike our intuition.
    For example, let $X$ be a variety over finite field $\bbF_q$ and consider the absolute Frobenius morphism $\Frob_q$.
    Then $\Frob_q$ is the identity morphism on the set of scheme-theoretic points of $X$.
    But on $X(\overline{\bbF_q})$, it has many interesting dynamical properties.
  \end{remark}

  Let $X,Y$ be a variety over $\kk$, and $f: X \dashrightarrow Y$ be a rational map.
  Recall that a rational map $f: X \dashrightarrow Y$ is a morphism defined on a dense open subset of $X$.
  The fundamental method to study a rational map is considering its graph.
  Let $\Gamma_f \subset X \times Y$ be the closure of the image of the rational map $(\id_X, f): X \dashrightarrow X \times Y$.
  Then $\Gamma_f$ is a variety birational to $X$ and admits two morphisms $\pi_1: \Gamma_f \to X$ and $\pi_2: \Gamma_f \to Y$ such that $f = \pi_2 \circ \pi_1^{-1}$.
  Using the graph, we can define the pullback and pushforward of divisors by $f$ as follows:
  \[ f^*D := \pi_{1*} \pi_2^* D, \quad f_*D := \pi_{2*} \pi_1^* D. \]
  However, not like morphisms, the pullback and pushforward by rational maps are not functorial in general, which causes some technical difficulties in the study of rational maps.
  We will fix this issue by considering $\bfb$-cycles on birational models in Section 2.2.

  \begin{example}
    Consider the rational map $f: \bbP^2 \dashrightarrow \bbP^2$ defined by $[x:y:z] \mapsto [1/x:1/y:1/z] = [yz:xz:xy]$.
    Then graph of $f$ is given by 
    \[ \Gamma_f = \{ ([x:y:z], [u:v:w]) \in \bbP^2 \times \bbP^2 : xu = yv = zw \}. \]
    Consider the first projection $\pi_1: \Gamma_f \to \bbP^2$ and $\bbA^2 \cong U \subset \bbP^2$ given by $z \neq 0$.
    Then 
    \begin{align*}
      \pi_1^{-1}(U) & = \{ ([x:y:1], [u:v:w]) \in \Gamma_f : xu = yv = w \} \\
      & = \{ ([x:y:1], [u:v]) \in U \times \bbP^1 : xu = yv \},
    \end{align*} 
    which is the blowup of $U$ at the origin.
    By similar argument, $\Gamma_f$ is the blowup of $\bbP^2$ at $[1:0:0]$, $[0:1:0]$ and $[0:0:1]$ and $\pi_1$ is the blowup morphism.

    We denote by $L_x$, $L_y$ and $L_z$ the coordinate lines in $\bbP^2$ given by $x=0$, $y=0$ and $z=0$ respectively, 
    and $E_x$, $E_y$ and $E_z$ the exceptional divisors in $\Gamma_f$ over $[1:0:0]$, $[0:1:0]$ and $[0:0:1]$ respectively.
    On $\Gamma_f$, the strict transform of $L_x$'s and the exceptional divisors $E_x$'s form a hexagon of rational curves.
    Taking $H = L_x$ on $\bbP^2$, then we have $p_2^* H = \tilde{L_x} + E_y + E_z$.
    When we push forward by $\pi_1$, the strict transform $\tilde{L_x}$ is contracted to a point, while $E_y$ and $E_z$ are mapped isomorphically to $L_y$ and $L_z$ respectively.
    Hence we have $f^* H = \pi_{1*} \pi_2^* H = L_y + L_z$.
    Then we get $(f^*)^2 H \sim f^*2H \sim 4H$ while $(f^2)^*H = (\id)^* H = H$.
    % \Yang{}
    % On $\tilde{X} = \Gamma_f$, the induced rational map $\tilde{f}$ is a morphism given by $([x:y:z], [u:v:w]) \mapsto ([u:v:w], [x:y:z])$.
    % The Neron-Severi group $\NS(\tilde{X})$ is given by
    % \[ \bigoplus_{i=x,y,z} \bbZ \tilde{L_i} \oplus \bbZ E_i \big/ \langle \tilde{L_i} + E_j + E_k \sim \tilde{L_{i'}} + E_{j'} + E_{k'} | \{i,j,k\} = \{i',j',k'\} = \{x,y,z\} \rangle. \]
    % Then we have $\tilde{f}^* \tilde{L_x} = \tilde{L_x}$, $\tilde{f}^* E_x = E_x$, $\tilde{f}^* E_y = \tilde{L_y}$ and $\tilde{f}^* E_z = \tilde{L_z}$.
  \end{example}

  We need to consider family of dominant rational maps.
  \begin{definition}
    Let $S$ be a noetherian integral scheme.
    A \emph{family of dominant rational maps} over $S$ is the following data:
    \begin{itemize}
      \item a scheme $X$ of finite type and separated over $S$ such that every fiber $X_s$ with $s \in S$ is a variety over $\kappa(s)$;
      \item an open dense subset $U \subseteq X$ and a morphism $f: U \to X$ over $S$ such that for every $s \in S$, the restriction of $f$ to the fiber $U_s$ gives a dominant rational map on $X_s$.
    \end{itemize}
    We denote by $f: X \dashrightarrow X$ the family of dominant rational maps over $S$.
  \end{definition}

  Given a family of dominant rational maps $f: X \dashrightarrow X$ over $S$, we can similarly define the graph $\Gamma_f \subset X \times_S X$ as the closure of the image of the morphism $(i_U, f): U \to X \times_S X$, where $i_U: U \to X$ is the inclusion morphism.
  We have $\Gamma_f$ is a scheme of finite type and separated over $S$ which is birational to $X$ via the first projection $\pi_1: \Gamma_f \to X$.

  \begin{proposition}
    Let $f: X \dashrightarrow X$ be a family of dominant rational maps over $S$.
    Then for every $s \in S$, $\Gamma_{f_s}$ is a closed subscheme of $(\Gamma_f)_s$.
    %  \Yang{}
  \end{proposition}
  \begin{proof}
    The commutative diagram 
    \[
      \begin{tikzcd}
        U_s \ar[r] \ar[d] & \Gamma_{f_s} \ar[r, hook] \ar[d,gray] & X_s \times_{\kappa(s)} X_s \ar[d] \\
        U \ar[r] & \Gamma_f \ar[r, hook] & X \times_S X
      \end{tikzcd}
    \]
    induces the gray arrow $\Gamma_{f_s} \to \Gamma_f$, and base changing by $s \to S$ gives the morphism $\Gamma_{f_s} \to (\Gamma_f)_s$.
    This morphism composing with the closed immersion $(\Gamma_f)_s \to X_s \times_{\kappa(s)} X_s$ yields the closed immersion $\Gamma_{f_s} \to X_s \times_{\kappa(s)} X_s$, so itself is a closed immersion.
  \end{proof}

  \begin{remark}
    In general, we have no equality $(\Gamma_f)_s = \Gamma_{f_s}$.
    For an example, see \Yang{}
  \end{remark}

  The following lemma is well-known.
  We will use it for lifting isolated periodic points from special fibers to the generic fiber.
  \begin{lemma}
    Let $X$ be a noetherian scheme and $Y \subset X$ be a regular immersion.
    For any subscheme $Z \subset X$ and $x \in Y \cap Z$ (if $Y \cap Z \neq \emptyset$), we have 
    \[ \dim_x (Y \cap Z) \geq \dim_x Y + \dim_x Z - \dim_x X. \]
  \end{lemma}
  \begin{proof}
    Suppose that $Y$ is locally defined by a regular sequence $f_1, \dots, f_r$ in the local ring $\calO_{X,x}$.
    Then we have
    \[ \dim_x (Y \cap Z) = \dim \calO_{Z,x} / (f_1, \dots, f_r) \geq \dim \calO_{Z,x} - r = \dim_x Z - (\dim_x X - \dim_x Y) \]
    by Krull's principal ideal theorem.
    % \Yang{}
  \end{proof}
  
  We may need the following special case of the flattening theorem of Gruson-Raynaud.
  \begin{theorem}[{Gruson-Raynaud, ref. \cite[Tag 0815]{Stacks}}]
    Let $S$ be a noetherian integral scheme and $X,Y$ of finite type and separated scheme over $S$.
    Let $f: X \to Y$ be a morphism over $S$ and $U \subseteq Y$ be an open subset such that $f^{-1}(U) \to U$ is flat.
    Then there exists a blowup $\pi: Y' \to Y$ with center in $Y \setminus U$ such that the induced morphism $f': X' \to Y'$ is flat, where $X'$ is the Zariski closure of $f^{-1}(U) \times_U \pi^{-1}(U)$ in $X \times_Y Y'$.
    %  \Yang{}
  \end{theorem}


\subsection{Intersection product on birational models}

  Let $X$ be a variety over $\kk$.
  All birational models of $X$ form a directed system, we denote by $\frakX$.
  Let $\CH^i(X)$ be the Chow group of cocycles of codimension $i$ on $X$; see \cite[Chapter 17]{Fulton13Intersection}.
  For our purpose, we only need the numerical equivalence classes of cocycles, namely $\upN^i(X)$; see \cite[Section 2]{Dang20degrees}.
  Since the pullback by rational maps is not functorial, we need to consider the inductive limit of the Chow groups $\CH^i(X')$ and $\upN^i(X')$ for all birational models $X' \in \frakX$.
  A dominant rational map $f: X \ratmap Y$ between varieties of the same dimension will induce a ``morphism'' $\frakX \to \frakY$ between the directed systems of birational models, which allows us to define a functorial pullback.
  
  \begin{definition}
    We define 
    \[ \CH^i(\frakX) := \varinjlim_{\pi:X'' \to X' \in \frakX} \left( \pi^* : \CH^i(X') \to \CH^i(X'') \right). \]
    An element in $\CH^i(\frakX)$ is called a \emph{$\bfb$-cocycle} of degree $i$ on $X$.
    It can be viewed as a pair $(X', \alpha' \in \CH^i(X'))$ modulo the equivalence relation generated by $(X_1, \alpha_1) \sim (X_2, \alpha_2)$ if there is a higher birational model $X'' \xrightarrow{\pi_i} X_i$ for $i=1,2$ such that $\pi_1^* \alpha_1 = \pi_2^* \alpha_2$.
    The model $X'$ is called a \emph{defining model} of $\alpha$.
    For $\alpha \in \CH^i(X)$, we still denote by $\alpha$ the $\bfb$-cocycle defined by $\alpha$.

    We can define $\upN^i(\frakX)$ similarly by taking the inductive limit of $\upN^i(X')$.
    %  \Yang{}
  \end{definition}

  The intersection product $\CH^i(X') \otimes \CH^j(X') \to \CH^{i+j}(X')$ induces an intersection product
  \[ \CH^i(\frakX) \otimes \CH^j(\frakX) \to \CH^{i+j}(\frakX) \]
  by taking the intersection product on a common defining model.
  It is well-defined since the pullback is a ring homomorphism.
  The intersection number 
  \[ \CH^i(\frakX) \otimes \CH^{d-i}(\frakX) \to \bbZ \]
  is also well-defined by the projection formula.

  Similarly, the intersection product and intersection number on $\upN^i(\frakX)$ are also well-defined.

  % On $\upN^i(X)$, there is a natural cone called the \emph{BPF} cone containing all intersections of nef divisors, denoted by $\Bpf^i(X)$; see \cite[Section 2]{Dang20degrees}.
  % The property of being in $\Bpf^i(X)$ is preserved by pullback, so we can also define $\Bpf^i(\frakX)$ as the inductive limit of $\Bpf^i(X')$ for all birational models $X' \in \frakX$.
  % One can view $\Bpf^i(\frakX)$ as a cone in $\upN^i(\frakX)$ whose elements are of form $(X', \alpha' \in \Bpf^i(X'))$ for some birational model $X' \in \frakX$.

  Let us focus on the case of divisors, i.e. $i=1$.
  In this case, note that the notion of nef, effective, big and pseudo-effective divisors are preserved by pullback, so we can also define these notions for $\bfb$-divisors on $\frakX$.
  In particular, although we cannot define ampleness for $\bfb$-divisors, we can still define the nef and big $\bfb$-divisors, which are enough for our purpose.

  Let $f: X \dashrightarrow Y$ be a dominant rational map between varieties of the same dimension.
  Then for every $\alpha \in \upN^i(\frakY)$, $f$ induces a rational map $X \ratmap Y'$ with $Y'$ a defining model of $\alpha$.
  Then $f$ induces a morphism $X' \to Y'$ for some birational model $X' \in \frakX$ by taking graph.
  The pullback $f^* \alpha$ is a well-defined element in $\upN^i(X')$, and hence defines an element in $\upN^i(\frakX)$.
  Easily check that the pullback $f^* \alpha$ is independent of the choice of the defining model $Y'$ of $\alpha$ and the choice of the birational model $X'$ of $X$.
  In this way, if $f: X \dashrightarrow Y$ and $g: Y \dashrightarrow Z$ are two dominant rational maps between varieties of the same dimension, then we have $(g \circ f)^* = f^* \circ g^*$ as homomorphisms $\upN^i(\frakZ) \to \upN^i(\frakX)$.


  


  \Yang{}
  % We define the inductive limit of the Picard groups $\Pic(X')$ by 
  % \[ \Pic(\frakX) \coloneqq \varinjlim_{\pi:X'' \to X' \in \frakX} \left( \pi^* \Pic(X') \to \Pic(X'') \right). \]
  % One can view $\Pic(\frakX)$ as the group of $\bfb$-Cartier divisors modulo linear equivalence on $X$.
  % For every $D \in \Pic(\frakX)$, there exists a birational model $X' \in \frakX$ and $D' \in \Pic(X')$ such that $D_{X''} = \pi^* D'$ for every $\pi: X'' \to X'$ in $\frakX$.
  % The model $X'$ is called a \emph{defining model} of $D$.

  % The intersection product on $\Pic(X')$ induces an intersection product on $\Pic(\frakX)$, which is well-defined by the projection formula.

  % Note that nefness or bigness of a divisor is preserved by pullback, so we can also define $L \in \Pic(\frakX)$ to be nef or big if $L_{X'}$ is nef or big for some (hence any) defining model $X' \in \frakX$.

  %  \Yang{}

  % \begin{remark}
  %   One can view $\frakX$ as the Riemann-Zariski space associated to $X$.
  %   But for our purpose, we only need to consider the inductive limit of the Picard groups, so we do not need to introduce the language of Riemann-Zariski space.
  % \end{remark}
  
  % \Yang{The intersection product on $\Pic(X')$ induces an intersection product on $\Pic(\frakX)$, which is well-defined by the projection formula.}


  % Let $L$ be a nef divisor on $X$.
  % Then the pullback of $L$ to every birational model of $X$ is also nef, so we can view $L$ as an element in $\Pic(\frakX)$.

  % \Yang{nef and big}


  \Yang{Intersect with curves}

  We also need to consider the intersection of $\bfb$-divisors with curves.
  Define
  \[ Z_1(\frakX) := \varprojlim_{\pi:X'' \to X' \in \frakX} \left( \pi_*: Z_1(X'') \to Z_1(X') \right), \]
  where the transition maps are given by pushforward of curves.

  For any curve $C$ on $X$, by Zorn's Lemma, there exists a $\bfb$-curve $\tilde{C}$ on $\frakX$ such that $\tilde{C}_X = C$.

  \Yang{functoriality of pushforward and pullback by rational maps}

  \begin{proposition}
    The projection formula holds for $\bfb$-divisors and $\bfb$-curves, i.e. for every $D \in \Pic(\frakX)$ and every $C \in Z_1(\frakX)$, and $f: \frakX \to \frakY$ be a morphism of Riemann-Zariski spaces, we have
    \[ D \cdot f_* C = f^* D \cdot C. \]
     \Yang{}
  \end{proposition}
  
  \begin{definition}
    The group $\upN^k(\frakX)$ and $\Bpf^k(\frakX)$ are defined as the inductive limit of $\upN^k(X')$ and $\Bpf^k(X')$ respectively, where the transition maps are given by pullback of divisors.
     \Yang{}
  \end{definition}

  \begin{theorem}[{Generalized Siu's inequality, ref. \cite{JiangLi2023Siu-inequality}}]
    Let $X$ be a projective variety of dimension $d$ and $\alpha,\beta,\eta$ be intersections of nef divisors on $X$ and $H$ a nef divisor on $X$.
    Then we have \Yang{}
    \[ (\alpha \cdot \beta) \cdot H^{d} \leq \binom{d}{k} (\alpha \cdot H^{d-k}) (\beta \cdot H^{k}). \]
  \end{theorem}

  \begin{corollary}
    Let $X$ be a projective variety of dimension $d$ and $\alpha,\beta,\eta$ be intersections of nef divisors on $X$ and $H$ a nef divisor on $X$.
    Then we have \Yang{}
    \[ (\alpha \cdot \beta \cdot \eta) \cdot (H^{k+l} \cdot \eta) \leq \binom{k+l}{k} (\alpha \cdot H^{l} \cdot \eta) (\beta \cdot H^{k} \cdot \eta). \]
  \end{corollary}

  \begin{theorem}[{Generalized Siu's inequality, ref. \cite{JiangLi2023Siu-inequality}}]
    Let $X$ be a projective variety of dimension $d$ and $\alpha \in \Bpf^k(X)$, $\beta \in \Bpf^{d-k}(X)$ and $H$ be a nef divisor on $X$.
    Then we have \Yang{}
    \[ (\alpha \cdot \beta) \cdot H^d \leq \binom{d}{k} (\alpha \cdot H^{d-k}) (\beta \cdot H^k). \]
  \end{theorem}

  \begin{theorem}[{ref. \Yang{}}]
    Let $\alpha \in \Bpf^k(\frakX)$ and $f: \frakX \to \frakY$ be a dominant rational map between spaces of the same dimension.
    Then $f_* \alpha \in \Bpf^k(\frakY)$.
  \end{theorem}





\subsection{Dynamical degrees}

  \begin{proposition}
    Let $f: X \dashrightarrow X$ be a rational map on a projective variety $X$ of dimension $d$.
    Then the limit 
    \[ \lambda_k(f) := \lim_{n \to \infty} \left(\deg_{k,H}(f^n)\right)^{1/n} \]
    exists and is independent of the choice of the ample divisor $H$.
    \Yang{}
  \end{proposition}

  \begin{proposition}
    The sequence of dynamical degrees $\lambda_k(f)$ is log-concave, i.e. 
    \[ \lambda_k(f)^2 \geq \lambda_{k-1}(f) \lambda_{k+1}(f)\] 
    for every $1 \leq k \leq d-1$.
     \Yang{}
    Or equivalently, the sequence of Lyapunov exponents $\mu_k(f)$ is non-increasing.
     \Yang{}
  \end{proposition}



  
\subsection{Semi-continuous functions on noetherian schemes}

  Let $S$ be an integral noetherian scheme, and $\theta: S \to \bbR$ be a function.
  We say that $f$ is \emph{lower semi-continuous} if for every $r \in \bbR$, the set $\{ s \in S : f(s) > r \}$ is open in $S$.

  \begin{remark}
    The limit of a sequence of lower semi-continuous functions may not be lower semi-continuous. 
    \Yang{}
  \end{remark}

  \begin{lemma}
    A function $\theta: S \to \bbR$ is lower semi-continuous if and only if the following condition holds: 
    \begin{enumerate}
      \item for every $\eta \in S$ and every specialization $\xi \in \overline{\{\eta\}}$, we have $\theta(\xi) \leq \theta(\eta)$;
      \item for every $\eta \in S$ and $a < \theta(\eta)$, there is an open subset $U \subseteq \overline{\{\eta\}}$ such that $\theta(\xi) > a$ for every $\xi \in U$.
    \end{enumerate}
  \end{lemma}
  \begin{proof}
    \Yang{}
  \end{proof}


  \begin{lemma}
    Let $\theta: S \to \bbR$ be a lower semi-continuous function.
    Then 
    \begin{enumerate}
      \item $\theta$ is continuous with respect to the constructible topology on $S$;
      \item if $\theta$ is discrete, then its image is finite.
    \end{enumerate}
  \end{lemma}
  \begin{proof}
    \Yang{}
  \end{proof}




\subsection{The Frobenius}

  In this subsection, we review some fundamental properties of the Frobenius morphism, which may be trivial for experts. 
  But for me, they are not so trivial at least at the beginning of this note.

  \begin{definition}
    Let $S$ be a scheme over $\bbF_p$.
    Then the \emph{absolute Frobenius} morphism $\Frob_{S,p}: S \to S$ is the morphism defined locally by the $p$-th power map on the structure sheaf $\calO_S$.
    %  \Yang{}
  \end{definition}

  \begin{proposition}
    Let $S,T$ be schemes over $\bbF_p$ and $f: S \to T$ be a morphism.
    Then $f$ is commutative with the absolute Frobenius $\Frob_p$.
  \end{proposition}

  \begin{definition}
    Let $S$ be a scheme over $\bbF_p$ and $X$ be a scheme over $S$.
    We denote by $X^{(p)}$ the base change of $X$ along the absolute Frobenius $\Frob_{S,p}$, i.e. $X^{(p/S)} = X \times_{S, \Frob_{S,p}} S$, called the \emph{Frobenius twist} of $X$.
    It is equipped with an $S$-scheme structure given by the second projection $X^{(p/S)} \to S$.

    The absolute Frobenius $\Frob_{S,p}: S \to S$ induces an $S$-morphism $\id_X \times \Frob_{S,p}: X \to X^{(p/S)}$, called the \emph{relative Frobenius} morphism of $X$ over $S$ and denoted by $\Frob_{X/S, p}$.
    In a diagram, we have
    \[ \begin{tikzcd}
      X \arrow[rrd, bend left, "\Frob_{X,p}"] \arrow[rdd, bend right] \arrow[rd, "\Frob_{X/S,p}"] & & \\
      & X^{(p/S)} \arrow[r] \arrow[d] & X \arrow[d] \\
      & S \arrow[r, "\Frob_{S,p}"] & S.
    \end{tikzcd} \]
    %  \Yang{}
  \end{definition}
  \begin{remark}
    We often omit the subscript $S$ on the Frobenius twist when there is no confusion, i.e. $X^{(p)}$ for $X^{(p/S)}$.
    We also often omit the subscript $X$ on the Frobenius morphism when there is no confusion, i.e. $\Frob_p$ for $\Frob_{X,p}$ and $\Frob_{/S,p}$ for $\Frob_{X/S,p}$.
    And we write $\Frob_{p^n}$ (resp. $\Frob_{/S,p^n}$) for the $n$-th iterate of $\Frob_p$ (resp. $\Frob_{/S,p}$).
  \end{remark}

  \begin{example}
    Let $S = \Spec \kkk$ be the spectrum of an algebraically closed field of characteristic $p$ and $X \subset \bbA^n_\kkk$ be an affine variety defined by $f_1, \dots, f_m \in \kkk[t_1, \dots, t_n]$.
    Then the Frobenius twist $X^{(p)}$ is given by $X^{(p)} = \Spec \kkk[t_1, \dots, t_n]/(f_1^{(p)}, \dots, f_m^{(p)})$, where $f_i^{(p)}$ is obtained by applying the $p$-th power map to the coefficients of $f_i$.
    The projection $X^{(p)} \to X$ is given by $(a_1, \dots, a_n) \mapsto (a_1^{1/p}, \dots, a_n^{1/p})$, while the relative Frobenius $\Frob_{/S,p}: X \to X^{(p)}$ is given by $(a_1, \dots, a_n) \mapsto (a_1^p, \dots, a_n^p)$.
    % \Yang{}
  \end{example}

  \paragraph{Relative Frobenius as endomorphisms when defined over finite fields}
  The most interesting case for us is when $S = \Spec \bbF_q$ or $S = \Spec \overline{\bbF_q}$ and $X$ is a variety over $S$.
  First consider the case when $S = \Spec \bbF_q$.
  Then the absolute Frobenius $\Frob_q: \Spec \bbF_q \to \Spec \bbF_q$ is the identity morphism, so the Frobenius twist $X^{(q)}$ is just $X$ and the relative Frobenius $\Frob_{/S,q}$ is just the absolute Frobenius $\Frob_q$.
  
  More generally, let $\kkk$ be an arbitrary algebraically closed field of characteristic $p$ and $X$ be a variety over $\kkk$.
  Suppose that $X$ is defined over some finite field $\bbF_q$.
  That is, there exists a variety $X_0$ over $\bbF_q$ and an isomorphism $X \cong X_0 \times_{\bbF_q} \kkk$.
  Then we have diagrams
  \[ 
    \begin{tikzcd}
      X \arrow[r] \arrow[rdd] \arrow[rd, red] & \Spec \kkk \arrow[rd, "\Frob_q"] \arrow[rdd] & \\
      & X \arrow[r] \arrow[d] & \Spec \kkk \arrow[d] \\
      & X_0 \arrow[r] & \Spec \bbF_q
    \end{tikzcd}
    \quad \leadsto \quad
    \begin{tikzcd}
      X \arrow[rrd, red] \arrow[rdd] \arrow[rd, blue] & & \\
      & X^{(q)} \arrow[r] \arrow[d] & X \arrow[d] \\
      & \Spec \kkk \arrow[r, "\Frob_q"] & \Spec \kkk.
    \end{tikzcd} 
  \]
  In particular, this gives an isomorphism $X^{(q)} \cong X$ (the blue arrow) over $\kkk$.
  In general, for a variety $X$ over $\kkk$, we still have an isomorphism $X^{(q)} \cong X$ as schemes but not as $\kkk$-schemes.
  Using this $\kkk$-scheme isomorphism, the relative Frobenius $\Frob_{/\kkk,q}: X \to X^{(q)}$ can be viewed as a morphism $X \to X$, which is different from the absolute Frobenius $\Frob_q: X \to X$.
  Another more direct way to get the relative Frobenius is to take the product 
  \[ \Frob_{/\kkk,q} = (\Frob_{q})_\kkk = \Frob_{X_0/\bbF_q, q} \times \id_\kkk: X_0 \times_{\bbF_q} \kkk \to X_0 \times_{\bbF_q} \kkk. \]
  This is the construction we will use in the following.

  If $\kkk = \overline{\bbF_q}$, then every variety over $\kkk$ is defined over some finite field, so the above construction applies to every variety over $\kkk$.
  
  \begin{proposition}
    Let $X$ be a variety over $\bbF_{q}$.
    Then for every closed point $x \in X$, $x \in X(\bbF_{q^n})$ if and only if $\Frob_{q}^n(x) = x$.
    In particular, every closed point of $X$ is a periodic point of the Frobenius $\Frob_q$.
    %  \Yang{}
  \end{proposition}
  \begin{proof}
    \Yang{}
  \end{proof}

  \begin{proposition}
    Let $f: X \to Y$ be a morphism of varieties over $\kkk$.
    Suppose that $f,X,Y$ are defined over some finite field $\bbF_q$.
    Then $f$ is commutative with the relative Frobenius $\Frob_{/\kkk,q}$.
    %  \Yang{}
  \end{proposition}
  \begin{proof}
    Choose $f_0: X_0 \to Y_0$ defined over $\bbF_q$ such that $f = f_0 \times_{\bbF_q} \kkk$.
    Since $f_0$ is commutative with the absolute Frobenius $\Frob_q$, the base change $f = f_0 \times_{\bbF_q} \kkk$ is also commutative with the relative Frobenius $\Frob_{/\kkk,q} = (\Frob_q)_\kkk$.
  \end{proof}
  


  % \Yang{}
  % \begin{example}
  %   Let $X = \bbA^1_\kk = \Spec \kk[t]$ be the affine line over $\kk$ and $X_\kkk = \Spec \kkk[t]$ be the base change of $X$ to $\kkk$.

  %   Then the Frobenius twist $X_\kkk^{(q)}$ is given by $X_\kkk^{(q)} = \Spec \kkk[t] \otimes_{\kkk, \Frob_q} \kkk$.
  %   We denote the ring $\kkk[t] \otimes_{\kkk, \Frob_q} \kkk$ by $R$ for simplicity.
  %   The elements in $R$ are of form $\sum_i a_i t^i \otimes b_i$ with $a_i, b_i \in \kkk$.
  %   The $\kkk$-algebra structure on $R$ is given by 
  %   \[ c (at^i \otimes b) = a t^i \otimes (c b) = (c^{1/q} a) t^i \otimes b \]
  %   for every $c \in \kkk$.
  %   There is a natural ring isomorphism $R \to \kkk[t]$ given by $\sum_i a_i t^i \otimes b_i \mapsto \sum_i a_i^q b_i t^i$.

  %   \Yang{}


  %   There is a natural ring isomorphism $R \to \kkk[t]$ given by $\sum_i a_i t^i \otimes b_i \mapsto \sum_i a_i^q b_i t^i$.
  %   However, use this isomorphism to identify $R$ with $\kkk[t]$, the $\kkk$-algebra structure on $\kkk[t]$ is given by $c \cdot t^i = c^q t^i$ for every $c \in \kkk$, which is different from the usual $\kkk$-algebra structure on $\kkk[t]$.
  %   We denote by this $\kkk$-algebra $\kkk[t]$ by $\kkk[t]^{(q)}$ for clarity.

  %   Now we can describe the absolute and relative Frobenius morphisms.
  %   The absolute Frobenius $\Frob_q: X \to X$ is given by 
  %   \[ \kkk[t] \to \kkk[t], \quad \sum a_i t^i \mapsto \sum a_i^q t^{qi}. \]
  %   It is not $\kkk$-linear, but it is $\bbF_q$-linear.
  %   The relative Frobenius $\Frob_{/S,q}: X \to X^{(q)}$ is also given by 
  %   \[ \kkk[t]^{(q)} \to \kkk[t], \quad \sum a_i t^i \mapsto \sum a_i^q t^{qi}. \]
  %   But this is $\kkk$-linear because of the $\kkk$-algebra structure on $\kkk[t]^{(q)}$.
    
  %   Since $\kkk$ is perfect, the absolute Frobenius $\Frob_q: \kkk \to \kkk$ is an isomorphism, so its base change 
  %   \[ \kkk[t] \to \kkk[t]^{(q)}, \quad \sum a_i t^i \mapsto \sum a_i^q t^i \]
  %   is also an isomorphism.
  %   Note that this isomorphism is $\kkk$-linear.
  %   Using this isomorphism to identify $\kkk[t]^{(q)}$ with $\kkk[t]$, the relative Frobenius becomes
  %   \[ \kkk[t] \to \kkk[t], \quad \sum a_i^q t^i \mapsto \sum a_i^q t^{qi}. \]

  %   On the other hand, the geometric Frobenius $\Frob_{X/\bbF_q, q}^{geo}: X_\kkk \to X_\kkk$ is given by
  %   \[ \kkk[t] \to \kkk[t], \quad \sum a_i t^i \mapsto \sum a_i t^{qi}. \]
    
  % \end{example}

  % \paragraph{Action on scheme-theoretic points and $\kk$-points.}

  %   \Yang{}
  %   Let $X$ be a scheme over $\bbF_p$ and $x \in X$ be a scheme-theoretic point.
  %   Then the absolute Frobenius $\Frob_p$ acts on $x$ by sending it to the point $\Frob_p(x)$ defined by the ideal sheaf $\calI_{\Frob_p(x)} = \calI_x^p$.
  %   In particular, if $x$ is a closed point, then $\Frob_p(x) = x$

  % \paragraph{Graph of the Frobenius.}
  % The graph of the absolute Frobenius $\Frob_{X,p}$ is given by


  % \paragraph{Absolute Frobenius as morphism over finite fields.}
  % When $\kk = \bbF_q$, the absolute Frobenius $\Frob_{q}$ is the identity morphism on $\Spec \bbF_q$.
  % Let $X$ be a variety over $\bbF_q$, then the Frobenius twist $X^{(q)}$ is isomorphic to $X$ and the absolute Frobenius $\Frob_{X, q}$ coincides with the absolute Frobenius $\Frob_{X/\bbF_q,q}$, which is a morphism over $\bbF_q$.
  % % so the Frobenius twist $X^{(q)}$ is isomorphic to $X$ and the absolute Frobenius $\Frob_{X, q}$ coincides with the absolute Frobenius $\Frob_{X/\bbF_q,q}$, which is a morphism over $\bbF_q$.

  % \begin{proposition}
  %   Let $X$ be a variety over $\overline{\bbF_p}$.
  %   Then every closed point $x \in X(\overline{\bbF_p})$ is a periodic point of the absolute Frobenius $\Frob_p$.
  %    \Yang{}
  % \end{proposition}

  % \begin{proposition}
  %   Let $X,Y$ be varieties over $\bbF_q$ and $f: X \to Y$ be a morphism over $\bbF_q$.
  %   Then $f$ is commutative with the absolute Frobenius $\Frob_q$.
  %    \Yang{}
  % \end{proposition}




